%- block legend
The jury has secretly chosen an integer $S$ such that $1 \leq S \leq N$. Your
task is to find $S$ using at most $\VAR{vars.Q.max}$ guesses.

This is a function-call problem. Implement the following function:

\begin{center}
\texttt{int find\_secret(int n);}
\end{center}

It receives $N$ ($\VAR{vars.N.min} \leq N \leq \VAR{vars.N.max}$) and must
return $S$. It is called exactly once. To make a guess, call:

\begin{center}
\texttt{int guess(int x);}
\end{center}

with $1 \leq x \leq N$. It returns $0$ if $x = S$, $1$ if $S > x$, and $-1$ if
$S < x$. You may call \texttt{guess} at most $\VAR{vars.Q.max}$ times.

Your submission must include \texttt{guessing.h} and must not define
\texttt{main} nor read from or write to the standard streams.
%- endblock

%- block interaction
The grader calls \texttt{find\_secret} and reports the value it returned. For
reference, the sample below shows the messages exchanged with the jury: the
grader sends $N$, then one line \texttt{? x} per call to \texttt{guess(x)},
followed by the result of that call, and finally \texttt{! s} once
\texttt{find\_secret} returns $s$.
%- endblock

%- block scoring
\begin{tabular}{|c|c|l|}
\hline
\textbf{Subtask} & \textbf{Points} & \textbf{Constraints} \\
\hline
%- for g in problem.groups
%- if g.score
\VAR{g.name | escape} & \VAR{g.score} & No additional constraints. \\
\hline
%- endif
%- endfor
\end{tabular}
%- endblock

%- block notes
No notes.
%- endblock
